\subsection{Kỹ năng mềm}

\subsubsection{Kỹ năng giao tiếp}
Trong quá trình triển khai dự án, giao tiếp rõ ràng và minh bạch là yếu tố quan trọng giúp nhóm duy trì nhịp độ làm việc:
\begin{itemize}
    \item \textbf{Trao đổi thường xuyên:} Tổ chức các cuộc họp ngắn định kỳ và trao đổi liên tục qua các kênh liên lạc để cập nhật tiến độ công việc cũng như các khó khăn kỹ thuật phát sinh.
    \item \textbf{Truyền đạt ý tưởng kỹ thuật:} Khi thiết kế các tính năng cho game Tetris (như cơ chế xoay khối, bảng điểm, tăng tốc độ rơi), các thành viên sử dụng tài liệu mô tả ngắn gọn, dễ hiểu để thống nhất phương án triển khai.
    \item \textbf{Phản hồi mang tính xây dựng:} Tiếp thu ý kiến đóng góp từ đồng đội trong các buổi đánh giá mã nguồn (code review) để cùng hoàn thiện sản phẩm.
\end{itemize}

\subsubsection{Kỹ năng làm việc nhóm}
Làm việc nhóm hiệu quả giúp tối ưu hóa thế mạnh của từng cá nhân và đảm bảo dự án hoàn thành đúng hạn:
\begin{itemize}
    \item \textbf{Phân chia công việc hợp lý:} Nhóm phân tích đề tài thành các module độc lập (logic game, giao diện console, viết tài liệu) và bàn giao cho từng thành viên phù hợp với năng lực.
    \item \textbf{Tinh thần tương trợ:} Sẵn sàng hỗ trợ, chia sẻ kiến thức khi thành viên khác gặp lỗi trong quá trình debug hoặc khó khăn khi thao tác với Git/LaTeX.
    \item \textbf{Đảm bảo cam kết tiến độ:} Mỗi thành viên có trách nhiệm hoàn thành đúng hạn các nhiệm vụ được giao để không gây ảnh hưởng đến tiến độ tổng thể của toàn nhóm.
\end{itemize}

\subsubsection{Kỹ năng giải quyết xung đột}
Quá trình phối hợp nhiều người không tránh khỏi những bất đồng về kỹ thuật và sự cố mã nguồn:
\begin{itemize}
    \item \textbf{Giải quyết bất đồng quan điểm:} Khi có sự khác biệt về thuật toán hoặc cách tổ chức code, nhóm tổ chức thảo luận, phân tích ưu nhược điểm dựa trên hiệu năng và tính tường minh để đưa ra quyết định chung.
    \item \textbf{Xử lý xung đột mã nguồn (Merge Conflict):} Khi xảy ra conflict lúc gộp nhánh trên GitHub, các thành viên liên quan trực tiếp rà soát từng dòng mã bị trùng lặp, thống nhất giữ lại logic đúng nhất trước khi hoàn tất merge.
\end{itemize}